\section{Some preliminaries}

\begin{definition}[Georgii (20.1)--(20.4)]
    \label{def:ir-20-1}
    \uses{def:rp-17-22}
    \lean{MeasureTheory.GibbsMeasure.heisenbergPotential, MeasureTheory.GibbsMeasure.sum_periodizedCoupling, MeasureTheory.GibbsMeasure.sum_abs_periodizedCoupling_le, MeasureTheory.GibbsMeasure.abs_heisenbergExponent_le, MeasureTheory.GibbsMeasure.HasExpSqMoments, MeasureTheory.GibbsMeasure.HasExpSqMoments.of_nonneg, MeasureTheory.GibbsMeasure.slowGrowth, MeasureTheory.GibbsMeasure.slowGrowth_sq, MeasureTheory.GibbsMeasure.IsSlowlySummable, MeasureTheory.GibbsMeasure.IsSlowlySummable.summable}
    \leanok

    The isotropic pair potential $\Phi_{\{i,j\}} = -J(i-j)\,\sigma_i\cdot\sigma_j$ of (17.22)
    is (20.1); $\|J\|_0 = \sum_k |J(k)| < \infty$ is (20.2), with the bound
    $|{}^\circ H_\Lambda| \leq \tfrac12\|J\|_0\sum_i|\sigma_i|^2$ (Georgii's constant is $\|J\|_0$);
    the a priori measure has exponential square moments (20.3), $\int e^{a|x|^2}\,d\lambda < \infty$
    for every $a$; and $J$ is slowly summable (20.4), $\sum_k |J(k)|\,\alpha(k)^2 < \infty$ with
    $\alpha(k) = (1 \vee \log|k|)^{1/2}$, which implies (20.2). Deviation: $|k|$ is the sup norm,
    which changes $\alpha$ by a bounded factor.
\end{definition}

\begin{lemma}[Georgii (20.6), on the torus]
    \label{lem:ir-20-6}
    \uses{def:ir-20-1, thm:rp-17-29, thm:rp-17-11}
    \lean{MeasureTheory.GibbsMeasure.measure_apply_pow_card_le, MeasureTheory.GibbsMeasure.measure_apply_pow_card_le_siteEquiv, MeasureTheory.GibbsMeasure.heisenbergPeriodicGibbsDist_apply_le_div, MeasureTheory.GibbsMeasure.heisenbergPeriodicGibbsDist_tail_le, MeasureTheory.GibbsMeasure.lintegral_exp_dotProduct_self_le_of_tail, MeasureTheory.GibbsMeasure.lintegral_indicator_abs_dotProduct_le, MeasureTheory.GibbsMeasure.slowGrowthSet, MeasureTheory.GibbsMeasure.measure_compl_slowGrowthSet_le, MeasureTheory.lintegral_exp_le_of_measure_le, summable_max_exp_norm_rpow, MeasureTheory.GibbsMeasure.integrable_exp_heisenbergExponent, MeasureTheory.GibbsMeasure.exists_lemma_20_6_of_tail, MeasureTheory.GibbsMeasure.heisenbergPeriodicField, MeasureTheory.GibbsMeasure.heisenbergPeriodicField_zero_tail_le, MeasureTheory.GibbsMeasure.exists_lemma_20_6_heisenbergPeriodicField, MeasureTheory.GibbsMeasure.heisenbergGZero, MeasureTheory.GibbsMeasure.exists_forall_tail_le_of_mem_heisenbergGZero, MeasureTheory.GibbsMeasure.exists_lemma_20_6_heisenbergGZero, MeasureTheory.GibbsMeasure.measurePreserving_shift_of_mem_heisenbergGZero, MeasureTheory.GibbsMeasure.sum_abs_periodizedCoupling_mul_slowGrowth_le, MeasureTheory.GibbsMeasure.map_periodicJuxt_apply_of_cylinderEvents, MeasureTheory.GibbsMeasure.measurePreserving_shift_map_periodicPullback}
    \leanok

    For a nonnegative-definite $J$ and an a priori measure with exponential square moments, the
    periodic Gibbs distribution satisfies the tail estimate
    ${}^\circ\gamma_\Lambda^\Phi(|\sigma_i|^2 \geq t) \leq b(a)e^{-at}$ uniformly in the cube,
    the site and the dimension, by the chessboard estimate (17.11) for a single-site event; hence,
    for any probability measure with this tail, (i) $\mu(e^{a|\sigma_i|^2}) \leq b_1(a)$, (ii)
    $\mu(|\sigma_i\cdot\sigma_j|\,1_{|\sigma_i\cdot\sigma_j| \geq t}) \leq b_2(a)e^{-at}$ (with a
    factor $2$ where Georgii uses Cauchy--Schwarz), and (iii) $\mu(\Omega_t^c) \leq b_3(a)e^{-at^2}$
    for the slow-growth set $\Omega_t = \{|\sigma_i|^2 \leq t^2\alpha(i)^2\ \forall i\}$, using
    $\sum_{\mathbb Z^d}(e \vee |i|)^{-(d+1)} < \infty$. Condition (20.3) makes
    ${}^\circ\gamma^\Phi_\Lambda$ well defined, as Georgii says in passing. The lemma as stated:
    for $\mu^\Phi_n = {}^\circ\gamma^\Phi_{\Lambda(n)} \times \delta_0$, the embedded periodic field
    with zero boundary condition (through (17.24), Theorem \ref{thm:rp-17-24}), the tail
    estimate holds at every site --- inside the cube through the coupling map, outside it because
    the spin is $0$ and $b(a) \geq 1$, rather than Georgii's ``by continuity'' --- and passes to
    the cluster points $\mathcal G_0(\Phi)$ of $(\mu^\Phi_n)$, so (i)--(iii) hold with three
    constants for every $\mu \in \{\mu^\Phi_n\} \cup \mathcal G_0(\Phi)$; every member of
    $\mathcal G_0(\Phi)$ is shift invariant, the part (d) of the proof of (20.7). Georgii's
    ``$J_{\Lambda(n)}$ is dominated using (20.4)'' is $\sum_{j\in\Lambda}|J_\Lambda(i-j)|\alpha(j)
    \leq \sum_k |J(i-k)|\alpha(k)$, uniformly in the cube.
\end{lemma}
\begin{proof}
    \leanok
\end{proof}

\begin{lemma}[Georgii (20.8)]
    \label{lem:ir-20-8}
    \uses{def:ir-20-1}
    \lean{MeasureTheory.exists_forall_le_dotProduct_of_measure_ne_zero, MeasureTheory.GibbsMeasure.heisenbergPeriodicGibbs_forall_mem_le, MeasureTheory.GibbsMeasure.le_heisenbergPeriodicGibbs_univ, MeasureTheory.GibbsMeasure.exists_le_lintegral_dotProduct_self_heisenbergPeriodicGibbsDist}
    \leanok

    There is $\delta > 0$, depending only on $\lambda$, with ${}^\circ\gamma_\Lambda^\Phi(|\sigma_i|^2)
    \geq \delta$ uniformly in the dimension, the cube, the site and $J$. Only $\lambda(\{0\}^c) > 0$
    is used, where Georgii assumes rotation invariance and non-degeneracy: a set $U$ with
    $\inf_{x,y\in U} x\cdot y > 0$ exists for any such measure. Two slips in his proof: the choice
    needs $\varepsilon^2 \leq u$, not $\varepsilon \leq u$; and the second estimate needs the lower
    bound $Z \geq \lambda(\{0\}\cup U)^{|\Lambda|}$ rather than the one he uses for the first.
\end{lemma}
\begin{proof}
    \leanok
\end{proof}

\begin{theorem}[Georgii (20.5): the Gibbsian specification of an unbounded N-vector model, and the gluing device of (13.18)]
    \label{thm:ir-20-5}
    \uses{def:ir-20-1, lem:ir-20-6}
    \lean{Specification.tailGlued, Specification.tailGlued_apply_of_mem, Specification.tailGlued_apply_of_notMem, Specification.measure_eq_one_of_isGibbsMeasure_tailGlued, Specification.IsPremodifierOn, Specification.piecewiseDensity, Specification.isPremodifier_piecewiseDensity, Specification.isSigmaFiniteLambdaAdmissible_piecewiseDensity, Specification.gibbsOnTailSet, Specification.gibbsOnTailSet_apply_of_mem, Specification.gibbsOnTailSet_apply_of_notMem, Specification.measure_eq_one_of_isGibbsMeasure_gibbsOnTailSet, Potential.IsSummableOn, Potential.hamiltonian_sub_eq_of_mem, Potential.tailBoltzmannFactor, Potential.isPremodifierOn_tailBoltzmannFactor, Potential.tailGibbsSpecification, Potential.tailGibbsSpecification_apply_of_mem, Potential.tailGibbsSpecification_apply_of_notMem, Potential.measure_eq_one_of_isGibbsMeasure_tailGibbsSpecification, MeasureTheory.GibbsMeasure.IsSlowlySummable.summable_sub, MeasureTheory.GibbsMeasure.slowGrowthTailSet, MeasureTheory.GibbsMeasure.mem_slowGrowthTailSet_iff, MeasureTheory.GibbsMeasure.measurableSet_tail_slowGrowthTailSet, MeasureTheory.GibbsMeasure.summable_abs_hamiltonianTerms_heisenbergPotential, MeasureTheory.GibbsMeasure.isSummableOn_heisenbergPotential, MeasureTheory.GibbsMeasure.exists_abs_hamiltonian_heisenbergPotential_le, MeasureTheory.GibbsMeasure.sigmaFiniteLambdaZ_tailBoltzmannFactor_ne_top, MeasureTheory.GibbsMeasure.heisenbergSpecification, MeasureTheory.GibbsMeasure.heisenbergSpecification_apply_of_mem, MeasureTheory.GibbsMeasure.heisenbergSpecification_apply_of_notMem, MeasureTheory.GibbsMeasure.heisenbergGibbsMeasures, MeasureTheory.GibbsMeasure.measure_slowGrowthTailSet_eq_one_of_mem_heisenbergGibbsMeasures, Matrix.abs_dotProduct_le_sqrt_mul_sqrt}
    \leanok

    The device of (13.18) and (20.5), once for all: a specification $\gamma$ glued along a tail
    event $\Omega_*$ with the Dirac specification $\delta_{0_\Lambda\omega_{S\setminus\Lambda}}$ off it is
    a specification (Georgii's (2.26) at a two-point parameter), and, as in (13.19), every Gibbs
    measure of the glued specification charges $\Omega_*$ fully whatever the branch on $\Omega_*$,
    provided $0 \in \Omega_*$. For a potential whose Hamiltonians converge only on $\Omega_*$, the
    Boltzmann factor is a premodifier on $\Omega_*$ (the relativised (1.31)), completed off
    $\Omega_*$ by $1$, and the Gibbsian branch is the resulting $\lambda$-specification, needing
    only $Z_\Lambda(\omega) < \infty$ on $\Omega_*$. (20.5) is the instance for the $N$-vector
    potential (20.1) over $\Omega_* = \bigcup_t \Omega_t$ of slowly increasing configurations, a
    tail event: (20.4) gives $\sum_j |J(i-j)|\alpha(j) < \infty$, Cauchy--Schwarz on the spin
    products gives the absolute convergence of $H^\Phi_\Lambda$ on $\Omega_*$ and Georgii's estimate
    $|H^\Phi_\Lambda(x)| \leq c(\omega_{S\setminus\Lambda})\sum_{i\in\Lambda}(1+|x_i|^2)$, and (20.3)
    gives $Z^\Phi_\Lambda(\omega) < \infty$; $\mu(\Omega_*) = 1$ for every $\mu \in \mathcal G(\Phi)$.
    The a priori measure need only be finite and non-zero; the completion by $1$ has no
    mathematical content, the gluing overwriting the kernel off $\Omega_*$. Not proved: (20.7)
    (the transfer of the tail estimate to the embedded periodic fields, the bound
    $\limsup_n\sup_{\Omega_t}\rho^n_\Lambda < \infty$, and the convergence $\gamma^n_\Lambda f \to
    \gamma^\Phi_\Lambda f$ on $\Omega_t$, since $\gamma^\Phi$ is not quasilocal and (4.17) as stated does
    not apply), and \S20.2.
\end{theorem}
\begin{proof}
    \leanok
\end{proof}

\begin{definition}[Georgii (20.9)--(20.10): the Fourier transform of the coupling and its restriction to the discrete dual]
    \label{def:ir-20-9}
    \uses{def:ir-20-1, def:gh-13-35}
    \lean{MeasureTheory.GibbsMeasure.couplingFourier, MeasureTheory.GibbsMeasure.ofReal_couplingFourier, MeasureTheory.GibbsMeasure.continuous_couplingFourier, MeasureTheory.GibbsMeasure.abs_couplingFourier_le, MeasureTheory.GibbsMeasure.couplingFourier_zero, MeasureTheory.GibbsMeasure.mFourier_zsmul_eq_one, MeasureTheory.GibbsMeasure.tsum_mul_mFourier_eq_sum_periodizedCoupling, MeasureTheory.GibbsMeasure.couplingFourier_eq_sum_latticeBox}
    \leanok

    $\hat J(z) = \sum_i z^i J(i)$ on the torus $G = K^d$ (Mathlib's characters of the unit torus,
    as in \S13.3), real for even $J$ and continuous for summable $J$, with $\hat J(1) = \sum_i J(i)$;
    on the $2n$-th roots of unity $G_\Lambda$ it is the finite sum $\sum_{i\in\Lambda} z^i
    J_\Lambda(i)$ of the periodised coupling (17.25). The correlation transform (20.11), the
    infrared bound (20.12)--(20.13) and Theorem (20.15) are not yet formalised: (20.13) is a
    three-stage approximation of the a priori measure combined with the chessboard estimate.
\end{definition}
